\section{Comparison of edge images}\label{sec:edge}
The local models of Section~\ref{sec:localization} now give an inclusion
between the two endpoint images of an edge. We prove only an inclusion
here, under freeness at the upper endpoint; Section~\ref{sec:induction}
will force equality and supply the remaining freeness.

Fix an interval isomorphism
\[
 \iota:[u,b]\longrightarrow [u',b']
\]
and a geometric reflection order in the second system.
By Theorem~\ref{bg:graph}, it orders the edges of the first interval
through their images under \(\iota\).
We call this the \emph{transported schedule}.
Order edges with equal labels arbitrarily.
Such edges share no vertex, since a reflection acts as an involution
without fixed vertices on its incident pairs.
In particular, the schedule totally orders the upward edges at each
fixed vertex.

Work with \(B(b)\) over \(A\). By Remark~\ref{bg:upper-interval}, every
upward edge at \(x\in[u,b]\) already lies in this interval, so the same
sets \(U_x\) and kernels \(L_x(F)\) are used here.

We first isolate the commutative-algebra fact used in the comparison.
\begin{lemma}\label{edge:lattice}
Let \(C\) be a polynomial ring over a field, let \(E\) be a finite free
\(C\)-module, and let \(M\subseteq E\) be a free submodule with the same
span over \(K=\Frac(C)\). Then, inside \(E\otimes_C K\),
\[
 M=\bigcap_{\substack{\mathfrak p\subset C\\
                     \operatorname{ht}\mathfrak p=1}}M_{\mathfrak p}.
\]
If \(C\) is a field, the same conclusion means \(M=E\).
\end{lemma}
\begin{proof}
Choose a \(C\)-basis of \(M\), which is also a \(K\)-basis of
\(E\otimes_C K\).
A vector in the intersection has coordinates in
\(\bigcap_{\operatorname{ht}\mathfrak p=1}C_{\mathfrak p}=C\).
For completeness, write a coordinate as a quotient of coprime
polynomials. Any nonconstant irreducible divisor of the denominator
defines a height-one prime at which the quotient does not belong
to the local ring. The denominator must therefore be a unit.
\end{proof}

\begin{proposition}[Edge-image bound]\label{edge:bound}
Let \(e:x\to y\) be an edge. Let \(F\subset U_x\) and \(G\subset U_y\)
consist of the upward edges strictly later than \(e\) in the transported
schedule. Put
\[
 C=A/(\alpha_e),\qquad
 I=\res_{x,e}L_x(F),\qquad J=\res_{y,e}L_y(G)
 \quad\text{inside }B^e .
\]
If \(L_y(G)\) is graded free over \(A\), there is a nonzero homogeneous
\(f_e\in C\) such that
\begin{equation}\label{eq:edge-bound}
 I\subseteq f_eJ,\qquad
 \deg f_e=\frac{d(x,y)-1}{2}.
\end{equation}
Moreover \(J\), and hence \(f_eJ\), is graded free over \(C\).
\end{proposition}
\begin{proof}
Write \(\alpha=\alpha_e\) and \(L=L_y(G)\).
For every \(g\in G\), the label \(\alpha_g\) is not proportional to
\(\alpha\): distinct reflections incident to a vertex have distinct
root lines. The module \(B^g\) is free over \(A/(\alpha_g)\).
Multiplication by \(\alpha\) is consequently injective on \(B^g\).
Applying this observation to all kernel conditions gives
\[
 L\cap\alpha B^y=\alpha L,\qquad J\simeq L/\alpha L.
\]
Indeed, if \(\alpha v\in L\), all restrictions of \(v\) along \(G\)
vanish by that injectivity.
The asserted freeness of \(J\) follows.

There is also no loss of generic rank.
For \(p_G=\prod_{g\in G}\alpha_g\), with the empty product equal to \(1\),
\[
 p_GB^y\subseteq L,\qquad
 \overline p_G B^e\subseteq J\subseteq B^e.
\]
The element \(\overline p_G\) is nonzero in \(C\).
Thus \(J\) spans \(B^e\) over \(\Frac(C)\).

For every root plane \(\Pi\) containing \(\alpha\), write \(D_\Pi\)
for its maximal dihedral subgroup and put
\[
 k_\Pi=\frac{\ell_{D_\Pi}(y)-\ell_{D_\Pi}(x)-1}{2}.
\]
These are nonnegative integers.
By Lemma~\ref{dih:sum}, only finitely many are nonzero and
\[
 \sum_\Pi k_\Pi=\frac{d(x,y)-1}{2}.
\]
Choose a root \(\beta_\Pi\in\Pi\) independent of \(\alpha\) whenever
\(k_\Pi>0\), and define
\[
 f_e=\prod_{k_\Pi>0}\overline{\beta_\Pi}^{\,k_\Pi}\in C .
\]
Different planes yield nonproportional reduced linear forms.
This is a nonzero homogeneous element of the required degree.

The target \(f_eJ\) is a free \(C\)-module of full generic rank in
\(B^e\). If \(C\) is a field, it equals \(B^e\), and the inclusion is
immediate. Otherwise Lemma~\ref{edge:lattice} reduces the inclusion to
height-one primes of \(C\).
At a prime containing no nonzero reduced root label, all factors in
\(\overline p_G\) and \(f_e\) are units.
Then \(J=B^e=f_eJ\) locally, and there is nothing to prove.
Every remaining height-one prime has the form
\(\mathfrak q=(\overline\beta)\), where \(\beta\) is a root independent
of \(\alpha\). Its inverse image in \(A\) is the plane prime
\(\mathfrak d=(\alpha,\beta)\). Put
\[
 \Pi=\operatorname{span}(\alpha,\beta),\qquad
 R=A_{\mathfrak d},\qquad
 \overline R=R/(\alpha)\cong C_{\mathfrak q}.
\]
The ring \(\overline R\) is a discrete valuation ring. Every root in
\(\Pi\) not proportional to \(\alpha\) reduces to a uniformizer times
a unit. Write \(I_{\mathfrak q}\) and \(J_{\mathfrak q}\) for the
localized images in \(B^e\otimes_C\overline R\).

Let \(K=D_\Pi x\cap[1,b]\), and choose the sheaf isomorphism
\[
 \psi:B_R|_K\longrightarrow
       \bigoplus_{c\in K}\mathcal O_{c,R}^{\oplus m_c}
\]
from Proposition~\ref{loc:decomposition}. The edge \(e\) is supported
in a model with top \(c\) exactly when \(y\le_{D_\Pi}c\). Thus the
edge component of \(\psi\) is an \(\overline R\)-linear isomorphism
\[
 \psi_e:B^e\otimes_C\overline R
   \longrightarrow
   \bigoplus_{\substack{c\in K\\y\le_{D_\Pi}c}}
            \overline R^{\oplus m_c}.
\]

The definitions of \(L_x(F)\) and \(L_y(G)\) are finite kernel
conditions, and their restrictions to \(e\) define \(I\) and \(J\)
as images. Flat localization commutes with these kernels and images.
The localized modules of edges outside \(\Pi\) are zero, so those
conditions disappear. Because \(\psi\) intertwines every restriction
map, its vertex components carry these localized kernels to the direct
sums of the corresponding model kernels. Its single edge component
\(\psi_e\) then carries both \(I_{\mathfrak q}\) and
\(J_{\mathfrak q}\) to the corresponding direct sums of model images.
Denote their ideals in a copy with top \(c\) by
\(I_c,J_c\subseteq\overline R\).

For this copy, the full relative interval \([x,c]_{D_\Pi}\) lies in
\([u,b]\): relative paths are ambient paths, with \(x\ge u\) and
\(c\le b\). By Remark~\ref{loc:arbitrary-edge-subsets}, the kernel at
either endpoint is generated by the product of the labels of its later
supported outgoing edges.

None of these labels is proportional to
\(\alpha\), so every factor reduces to a uniformizer times a unit.
The valuations of \(I_c\) and \(J_c\) are therefore the later-edge
counts \(h_x\) and \(h_y\) in \([x,c]_{D_\Pi}\).
Lemma~\ref{dih:count} gives \(h_x-h_y=k_\Pi\).
At \(\mathfrak q\), the element \(f_e\) has valuation \(k_\Pi\): all
other plane contributions are units, and the contribution from \(\Pi\)
is the unit \(1\) when \(k_\Pi=0\). Hence \(I_c=f_eJ_c\).

This equality holds in every supported copy. Since \(\psi_e\) is
\(\overline R\)-linear, it gives
\[
 \psi_e(I_{\mathfrak q})
 =f_e\psi_e(J_{\mathfrak q})
 =\psi_e(f_eJ_{\mathfrak q}),
\]
and thus \(I_{\mathfrak q}=f_eJ_{\mathfrak q}\) as embedded submodules
of the original localized edge module. Only \(\overline R\)-linearity is
used here, so the ungraded sheaf decomposition suffices.

We have proved \(I_{\mathfrak q}\subseteq f_eJ_{\mathfrak q}\) at every
height-one prime.
Lemma~\ref{edge:lattice} gives \(I\subseteq f_eJ\) globally.
The original maps and \(f_e\) are homogeneous, so the inclusion is graded.
\end{proof}

\begin{remark}\label{edge:codimension}
The argument tests membership in the free \emph{target} \(f_eJ\).
It makes no reflexivity assumption on \(I\), and does not infer
global equality from equality at height one.
A quotient \((f_eJ)/I\) supported in higher codimension could still
exist at this stage. It will be eliminated by the equality argument
below.
\end{remark}

\section{Reciprocity and the simultaneous induction}\label{sec:induction}
We prove Theorem~\ref{thm:main} together with the following statement:
\begin{quote}
For every interval isomorphism and every transported schedule,
the modules \(L_u(F)\) are graded free over \(A\) for all terminal
portions \(F\) of the upward edges at the bottom vertex \(u\).
\end{quote}
The statement includes the empty terminal portion and the full set.
It is quantified over both systems, so we may also interchange them.

Induct on \(d(u,b)\).
The statement is quantified over all choices of ambient systems and
realizations as above. Within each comparison, keep those data and the sheaf
\(B(b)\) fixed: passage to \([y,b]\) changes only the bottom vertex.
For rank zero the polynomial is \(1\), there are no edges, and the
only kernel is the free top stalk.
Suppose now that \(d(u,b)>0\).
For every \(y>u\), restriction of \(\iota\) gives an isomorphism
\([y,b]\cong[\iota(y),b']\) of smaller rank.
Restriction of the same reflection order gives exactly the schedule
on its edges.
The induction hypothesis therefore supplies both polynomial equality
at \(y\) and freeness for every terminal upward kernel at \(y\).
In particular, the hypothesis of Proposition~\ref{edge:bound}
is available at the upper endpoint of every edge in \([u,b]\).
No partial-kernel freeness at \(u\) is assumed.

Figure~\ref{fig:comparison-dependencies} records how these two induction
hypotheses lead to the bottom assertions.
\begin{figure}[!ht]
  \centering
  \begin{tikzpicture}[
    node distance=5mm and 7mm,
    comparisonstep/.style={
      draw=black!65,
      rounded corners=1.5pt,
      line width=.45pt,
      text width=5.7cm,
      align=center,
      inner xsep=5pt,
      inner ysep=4pt,
      font=\small
    },
    comparisonwide/.style={comparisonstep,text width=12.45cm},
    dependency/.style={-{Stealth[length=2mm,width=1.4mm]},line width=.55pt}
  ]
    \node[comparisonstep,minimum height=12mm] (upperfree) {
      \textbf{Upper terminal kernels are free}\\
      at every $y>u$, in both systems.
    };
    \node[comparisonstep,minimum height=12mm,right=of upperfree] (upperequal) {
      \textbf{Upper polynomials agree}\\
      $P_{y,b}=P'_{\iota(y),b'}$ for every $y>u$.
    };
    \node[comparisonstep,minimum height=16mm,below=of upperfree] (comparison) {
      \textbf{Hilbert comparison}\\
      The inclusion $I\subseteq f_eJ$ bounds each increase in Hilbert series.
    };
    \node[comparisonstep,minimum height=16mm,below=of upperequal] (reciprocity) {
      \textbf{Reciprocity and cancellation}\\
      Cancel the upper terms; compare in both directions.
    };
    \path (comparison.south west) -- (reciprocity.south east)
      coordinate[midway] (middle);
    \node[comparisonwide,below=6mm of middle] (polynomial) {
      \textbf{Polynomial equality at the bottom}\\
      The strict degree bound gives $P_{u,b}=P'_{u',b'}$.
    };
    \node[comparisonwide,below=of polynomial] (images) {
      \textbf{Equality of all edge images}\\
      Exact endpoint comparison forces every bound to be attained: $I=f_eJ$.
    };
    \node[comparisonwide,below=of images] (freeness) {
      \textbf{Bottom terminal kernels are free}\\
      Reverse the bottom schedule, starting from $L_u(\emptyset)=B^u$.
    };
    \draw[dependency] (upperfree) -- (comparison);
    \draw[dependency] (upperequal) -- (reciprocity);
    \draw[dependency] (comparison) -- (reciprocity);
    \draw[dependency] (reciprocity.south) -- (reciprocity.south |- polynomial.north);
    \draw[dependency] (polynomial) -- (images);
    \draw[dependency] (images) -- (freeness);
  \end{tikzpicture}
  \caption{The two uses of smaller-rank induction. Upper freeness supplies
    the edge bound; upper polynomial equality permits row cancellation
    against reciprocity. Polynomial equality then forces the edge bounds
    to be exact, and the reverse bottom schedule proves terminal-kernel
    freeness.}
  \label{fig:comparison-dependencies}
\end{figure}


\subsection{A nonnegative comparison}
Fix a real number \(q\) with \(0<q<1\), and put
\[
 T=q^{-1/2}-q^{1/2}>0.
\]
Recall that \(A\) has \(N\) polynomial variables.
For each \(x\in[u,b]\), set
\[
 p_x(q)=q^{-d(x,b)/2}P_{x,b}(q),\qquad
 p'_x(q)=q^{-d(x,b)/2}P'_{\iota(x),b'}(q).
\]
Regard $p=(p_x)_{x\in[u,b]}$ and $p'=(p'_x)_{x\in[u,b]}$ as
column vectors indexed by the first interval. The rank differences agree
because $\iota$ is a poset isomorphism.

Process all interval edges in increasing schedule order.
At any time, let \(H_x\) be the set of unprocessed upward edges at
\(x\), and keep the vector
\begin{equation}\label{eq:state}
 D_x=q^{-d(x,b)/2}(1-q)^N\Hilb(L_x(H_x);q).
\end{equation}
Each kernel is finitely generated and bounded below in degree,
so this Hilbert series converges at the chosen \(q\), regardless
of whether the kernel is free.
Proposition~\ref{bg:normalization} gives the endpoints
\begin{equation}\label{eq:endpoints}
 D_{\mathrm{initial}}=p(q^{-1}),\qquad
 D_{\mathrm{final}}=p(q).
\end{equation}

At \(e:x\to y\), removing its kernel condition gives the exact sequence
\[
 0\longrightarrow L_x(F\cup\{e\})
 \longrightarrow L_x(F)\longrightarrow I\longrightarrow0,
\]
where \(F\) is the strictly later set.
The current set at \(y\) is \(G\) from Proposition~\ref{edge:bound};
edges tied with \(e\) do not meet \(y\).
Only coordinate \(x\) changes.
Since \(J=L_y(G)/\alpha_e L_y(G)\),
\[
 \Hilb(f_eJ;q)
 =q^{(d(x,y)-1)/2}(1-q)\Hilb(L_y(G);q).
\]
The increase of \(D_x\) is consequently at most
\begin{align}
 &q^{-d(x,b)/2}(1-q)^N\Hilb(f_eJ;q)\notag\\
 &\quad=
 q^{-d(x,b)/2}(1-q)^{N+1}
 q^{(d(x,y)-1)/2}\Hilb(L_y(G);q)
 =T D_y. \label{eq:step}
\end{align}

Let \(E_{xy}\) be the matrix unit whose only nonzero entry is \(1\)
in row \(x\), column \(y\), and associate
\(M_e=\id+T E_{xy}\) to \(e\).
All entries of these matrices are nonnegative.
Repeatedly applying \eqref{eq:step} bounds the final vector by
the corresponding matrix product applied to the initial vector.

This product counts \emph{decreasing} reflection-label paths.
Indeed, an upper edge \(y\to z\) must already have transferred
its contribution into coordinate \(y\) before a lower edge
\(x\to y\) can use it.
For a whole path the edge labels therefore decrease from \(x\)
toward its final vertex.
Conversely, each such path appears once in the matrix expansion.
Equal labels cannot occur in a composable pair, and disjoint tied
edges have commuting matrices.
Define the matrix
$\widetilde R'(T)=(\widetilde R'_{\iota(x),\iota(z)}(T))_{x,z\in[u,b]}$,
with rows and columns indexed through $\iota$.
Theorem~\ref{bg:paths}, applied in the second system, gives
\begin{equation}\label{eq:comparison}
 p(q)\le \widetilde R'(T)p(q^{-1})
 \quad\text{entrywise},
\end{equation}


By the defining reciprocity identity in the second system,
evaluated at \(q^{-1}\), and by
\[
 (q^{-1})^{-d(x,z)/2}R'_{\iota(x),\iota(z)}(q^{-1})
   =\widetilde R'_{\iota(x),\iota(z)}(T),
\]
we have the exact identity
\begin{equation}\label{eq:reciprocal-vector}
 p'(q)=\widetilde R'(T)p'(q^{-1}).
\end{equation}
The positive sign of \(T\) here comes from evaluation at \(q^{-1}\).

\subsection{Polynomial equality}
The induction hypothesis gives \(p_y=p'_y\) for all \(y>u\).
Subtracting \eqref{eq:reciprocal-vector} from
\eqref{eq:comparison} in row \(u\), and using
\((\widetilde R'(T))_{u,u}=1\), yields
\[
 \Delta(q)\le\Delta(q^{-1}),\qquad \Delta=p_u-p'_u .
\]
Interchanging the two systems gives the reverse inequality.
Both comparisons use only smaller-rank freeness at upper endpoints.
Thus
\[
 \Delta(q)=\Delta(q^{-1})\qquad(0<q<1).
\]
The strict degree bound on Kazhdan--Lusztig polynomials implies that
\(\Delta\), as a Laurent polynomial in \(q^{1/2}\), has only strictly
negative exponents. Its reciprocal has only strictly positive
exponents. Equality on the interval \((0,1)\) forces both to vanish.
Hence \(P_{u,b}=P'_{u',b'}\).

\subsection{Equality forces all edge images}
It remains to prove the supplementary freeness statement.
The polynomial equality just obtained, together with the proper
upper equalities, makes the entire endpoint comparison
\eqref{eq:comparison} an equality.

Write the edge schedule as \(e_1,\ldots,e_r\), let \(D_k\) be the
actual state after step \(k\), and put \(M_k=M_{e_k}\).
Define the nonnegative deficit vector
\[
 \delta_k=M_kD_{k-1}-D_k\ge0.
\]
An exact telescoping identity gives
\begin{equation}\label{eq:deficit}
 M_r\cdots M_1D_0-D_r
 =\sum_{k=1}^r M_r\cdots M_{k+1}\delta_k.
\end{equation}
Every factor has nonnegative entries and diagonal entries \(1\);
so does every product, since each off-diagonal entry goes from a
strictly higher vertex to a lower one in the Bruhat order.
In particular, a positive coordinate of \(\delta_k\) remains positive
in the corresponding summand on the right.
The left side is zero by endpoint equality.
Therefore every deficit is zero.

At a particular edge, its only possibly nonzero coordinate is a
positive scalar times
\[
 \Hilb\big((f_eJ)/I;q\big).
\]
A nonzero finitely generated graded module has a strictly positive
Hilbert series at \(0<q<1\), including if its support has codimension
at least two or it has finite length.
Consequently the vanishing of the deficit proves
\begin{equation}\label{eq:image-equality}
 I=f_eJ
\end{equation}
as graded modules globally.
This supplies the equality deliberately not inferred from the
codimension-one tests.

\subsection{Freeness closes the induction}
Each image in \eqref{eq:image-equality} is graded free over
\(A/(\alpha_e)\); it therefore has a graded free \(A\)-resolution
of length at most one.
At the bottom \(u\), work backward through its upward-edge schedule,
starting with \(L_u(\varnothing)=B^u\).
Adding back one condition gives
\[
 0\longrightarrow K\longrightarrow F\longrightarrow I\longrightarrow0,
\]
where \(F\) is the previously reached free kernel and
\(\pd_A I\le1\).
Comparing this sequence with a length-one graded free resolution
of \(I\), graded Schanuel's lemma makes \(K\) graded projective.

A finitely generated graded projective module over a positively
graded polynomial ring whose degree-zero part is a field is graded
free. Here is the relevant argument.
Lift a homogeneous basis of \(K/A_{>0}K\) to obtain a graded
surjection from a finite graded free module \(F_0\) onto \(K\).
Graded projectivity splits it.
The complementary kernel has zero quotient modulo \(A_{>0}\),
and graded Nakayama makes that kernel zero.
Thus \(F_0\simeq K\).

This reverse procedure reaches every terminal portion of \(U_u\).
It proves the supplementary freeness assertion and closes the
simultaneous induction, establishing Theorem~\ref{thm:main}.
\qed

\begin{corollary}\label{cor:all-subintervals}
Every poset isomorphism \(\iota:[u,b]\to[u',b']\) preserves all
Kazhdan--Lusztig polynomials on its subintervals:
\[
 P_{x,y}(q)=P'_{\iota(x),\iota(y)}(q)
 \qquad(u\le x\le y\le b).
\]
It also preserves the corresponding \(R\)-polynomials.
\end{corollary}
\begin{proof}
Apply Theorem~\ref{thm:main} to each restricted interval.
For the second assertion, write reciprocity as
\[
 q^{d(x,y)}P_{x,y}(q^{-1})
 =\sum_{x\le z\le y}R_{x,z}(q)P_{z,y}(q).
\]
The term with \(z=y\) is \(R_{x,y}\), and all other terms use smaller
rank differences for \(R\).
Induction therefore recovers all \(R\)-polynomials from the
order and the preserved \(P\)-polynomials.
\end{proof}

\begin{remark}[Ranks of terminal kernels]\label{rem:terminal-ranks}
Use the chosen reflection order in \(W'\) also for its native schedule on
\([u',b']\), with the same order for corresponding tied edges in the two
global schedules. Let \(L'_{\iota(x)}\) denote the kernels of \(B'(b')\)
over its polynomial ring \(A'\), with \(N'\) variables. For each
\(x\in[u,b]\) and terminal
portion \(F\) at \(x\), its image \(F'=\iota(F)\) is a native terminal
portion at \(\iota(x)\), and
\[
 \operatorname{grk}_A(L_x(F);q)
 =\operatorname{grk}_{A'}(L'_{\iota(x)}(F');q).
\]
Indeed, \eqref{eq:image-equality} makes every transition exact, so
\(D_k=M_k\cdots M_1p(q^{-1})\). Apply the same argument to the identity
comparison in the second system. Corollary~\ref{cor:all-subintervals}
identifies the initial vectors, and the matched schedules give the same
matrices under \(\iota\). For each chosen pair \(x,F\), the two state coordinates therefore
agree at a common stage realizing \(F,F'\); such a stage exists because
incident edges are untied. The auxiliary assertion on \([x,b]\) and on
the identity comparison of \([\iota(x),b']\) makes both kernels free.
Their separate factors \((1-q)^N\) and \((1-q)^{N'}\) in the two versions of \eqref{eq:state}
then convert their Hilbert series to graded ranks, and the common length
factor cancels. Equality for all \(0<q<1\) proves the polynomial identity.
\end{remark}
